% !TEX root = ../main.tex
\documentclass[../main.tex]{subfiles}

\begin{document}

\selectlanguage{vietnamese}

\section{CHƯƠNG 3: BIỂU DIỄN TRI THỨC VÀ SUY LUẬN (Knowledge Representation and Reasoning)}

\subsection{3.1. Hệ thống dựa trên tri thức (Knowledge based systems)}
Một hệ thống dựa trên tri thức (\textit{Knowledge-based system}) tổ chức mã nguồn AI thành hai phần độc lập: Cơ sở tri thức (\textit{Knowledge Base - KB}) chứa tập hợp câu mệnh đề biểu diễn các sự thật về thế giới thực bằng một ngôn ngữ trang trọng, và Công cụ suy luận (\textit{Inference Engine}) chứa các thuật toán suy diễn logic độc lập với nội dung tri thức. Đại lý tương tác với hệ thống qua vòng lặp:
\begin{equation}
    \text{Tri giác mới} \xrightarrow{} \texttt{TELL}(KB) \xrightarrow{} \text{Công cụ suy luận xử lý} \xrightarrow{} \texttt{ASK}(KB) \rightarrow \text{Hành động tối ưu}
\end{equation}

\subsection{3.2. Logic mệnh đề (Propositional logic)}
\begin{itemize}
    \item \textbf{Cú pháp và Ngữ nghĩa:} Biểu diễn tri thức bằng các biến mệnh đề ký hiệu ($P, Q, R$) và các liên kết toán học ($\wedge, \vee, \neg, \Rightarrow, \Leftrightarrow$). 
    \item \textbf{Mô hình hóa toán học của Luật kéo theo logic (Entailment):} Một mô hình $m$ là một phép gán giá trị Đúng/Sai cho tất cả các biến logic. Gọi $M(KB)$ là tập hợp tất cả các mô hình làm cho cơ sở tri thức đúng, và $M(\alpha)$ là tập các mô hình làm cho câu mệnh đề $\alpha$ đúng. Biểu thức toán học thiết lập đặc trưng của hệ suy luận:
    \begin{equation}
        KB \models \alpha \iff M(KB) \subseteq M(\alpha)
    \end{equation}
    Nghĩa là câu $\alpha$ là hệ quả logic của $KB$ nếu và chỉ nếu $\alpha$ luôn đúng trong tất cả các thế giới mà $KB$ đúng.
\end{itemize}

\subsection{3.3. Logic bậc nhất (First-order logic)}
Mở rộng logic mệnh đề để biểu diễn tri thức một cách tự nhiên và gọn nhẹ hơn nhờ việc đưa vào cấu trúc toán học của: Đối tượng (\textit{Objects} - thực thể trong thế giới), Thuộc tính và Mối quan hệ (\textit{Relations/Predicates}), Hàm (\textit{Functions}), và hai loại Lượng từ toán học ràng buộc biến:
\begin{itemize}
    \item \textbf{Lượng từ toàn thể ($\forall$):} Khẳng định tính chất đúng với mọi đối tượng trong miền xác định, thường đi kèm liên kết kéo theo ($\Rightarrow$).
    \item \textbf{Lượng từ tồn tại ($\exists$):} Khẳng định tính chất đúng với ít nhất một đối tượng trong miền xác định, thường đi kèm liên kết hội ($\wedge$).
\end{itemize}

\subsection{3.4. Suy luận trong Logic bậc nhất (Inference using FOL)}
Cơ chế suy luận dịch chuyển từ việc lập bảng chân trị bùng nổ không gian trạng thái sang các phép biến đổi trực tiếp trên câu tri thức: Phép thế (\textit{Substitution} $\{x/k\}$), Phép đồng nhất hóa (\textit{Unification}) tìm phép thế $\theta$ làm cho hai vị từ phức tạp trở nên giống nhau: $Unify(P(x), P(A)) = \{x/A\}$. Áp dụng các luật suy luận vạn năng như Modus Ponens mở rộng (\textit{Generalized Modus Ponens}).

\subsection{3.5. Thuật toán phân giải bác bỏ (Resolution refutation)}
Thuật toán phân giải bác bỏ là nền tảng của các công cụ suy luận tự động tối ưu nhất, hoạt động dựa trên phương pháp chứng minh phản chứng.
\begin{itemize}
    \item \textbf{Biểu thức toán học thiết lập đặc trưng hệ thống:} Để chứng minh câu $\alpha$ là đúng từ cơ sở tri thức, thuật toán thêm phủ định $\neg\alpha$ vào $KB$ và tiến hành triệt tiêu các bổ đề đối ngẫu để tìm ra mệnh đề rỗng đại diện cho mâu thuẫn toán học ($\emptyset$):
    \begin{equation}
        KB \wedge \neg\alpha \models \emptyset
    \end{equation}
    \item \textbf{Quy tắc phân giải tổng quát (Resolution Rule):} Yêu cầu chuyển đổi tất cả các câu tri thức phức tạp về Dạng chuẩn hội (\textit{CNF - Conjunctive Normal Form} - chuỗi các phép hội giữa các mệnh đề disjunction). Biểu thức khử đại số đặc trưng:
    \begin{equation}
        \frac{l_1 \vee \dots \vee l_i \vee \dots \vee l_k, \quad m_1 \vee \dots \vee m_j \vee \dots \vee m_n}{(l_1 \vee \dots \vee l_{i-1} \vee l_{i+1} \vee \dots \vee l_k \vee m_1 \vee \dots \vee m_{j-1} \vee m_{j+1} \vee \dots \vee m_n)\theta}
    \end{equation}
    Trong đó, $l_i$ và $m_j$ là hai bổ đề đối ngẫu có khả năng đồng nhất hóa thông qua phép thế $\theta$ sao cho $l_i\theta = \neg (m_j\theta)$.
\end{itemize}

\subsection{3.6. Các hệ thống suy luận logic (Logic reasoning systems)}
Bao gồm các cấu trúc lập trình lập luận tiến (\textit{Forward Chaining}) kích hoạt từ dữ liệu thô đẩy lên mục tiêu, lập luận lùi (\textit{Backward Chaining}) kích hoạt từ giả thuyết đích hạ xuống điều kiện cần, và các hệ thống ngôn ngữ lập trình logic như Prolog.

\subsection{3.7. Bài tập (Exercises)}
Thực hiện các bài toán chuyển đổi ngôn ngữ tự nhiên sang câu FOL, lập quy trình phân giải CNF để chứng minh các định lý toán học logic.

\end{document}
